% !TEX root = ../main.tex
\section{Discussion}

\hypertarget{what-drives-the-inconsistency}{%
\subsection{What drives the inconsistency}\label{what-drives-the-inconsistency}}

The reproduction cleanly separates the ingredients of the phenomenon. The harvest-flow constraint is the operative between-period link: without it, the multi-period problem decomposes and open-loop and sequential solutions agree; with it, and the tighter it is, the larger the divergence between plan and realization. The disequilibrium forest structure is the fuel: a mature/old-growth landbase places a large pulse of immediately harvestable, financially over-mature volume in front of a regulated future that cannot sustain it, so the open-loop plan's promised smooth flow can only be constructed by commitments the future planner will not keep. The negatively valued strata are the tell: because CM-CE area enters the solution only to mitigate the binding flow constraint and is then abandoned under replanning, its presence in the open-loop solution is a direct, observable marker of the inconsistency rather than a real management intention.

\hypertarget{the-discount-rate-nuance}{%
\subsection{The discount-rate nuance}\label{the-discount-rate-nuance}}

One finding deserves emphasis because it is easy to misstate, and misstatements of it circulate. The naive account---``a discounted objective plus an even-flow constraint breaks Bellman's principle''---suggests that the discount rate is the operative driver and that consistency could be restored by lowering it. The reproduction shows this is not so: occurrence and magnitude are identical at 0\%, 2\%, 4\%, and 6\%. The resolution is in the thesis itself (Daugherty 1991, ch.~3, p.~58--59): because the discount factor is uniform across periods, it cancels out of the conditions that a plan must satisfy to be consistent. Discounting shapes \emph{which} plan is chosen, but the inconsistency is a property of the constraint structure and the forest state, not of the rate. We present this as a reproduced nuance of the thesis, not as a contradiction of it---indeed, reproducing this counter-intuitive invariance on independent software is among the strongest checks that the reproduction has captured the thesis's mechanism rather than a look-alike artifact. The practical corollary is uncomfortable: discount-rate policy is not a remedy for non-credible plans.

\hypertarget{implications-for-the-literature}{%
\subsection{Implications for the literature}\label{implications-for-the-literature}}

The implication for the large literature of open-loop LP forest-planning studies---and for planning practice built on single-solve strategic models---is direct. Where a model combines a harvest-flow constraint with a disequilibrium initial forest, its projected harvest trajectories, residual growing-stock trajectories, and shadow prices are properties of a plan that sequential re-optimization will not follow, and trade-off analysis conducted on them is biased. Gunn's \citep[pp.~331--332]{gunn2007models} ``declining non-declining yield'' names the symptom exactly: the promised non-declining flow declines precisely because each successive planner declines to honour it. The same mechanism surfaces downstream, in the hierarchical link between strategic and operational planning: \citet{paradis2013risk} showed that incoherent linkage between long-term wood-supply allocation and short-term harvest planning induces a systematic drift between planned and implemented harvest (a divergence of forest system state), citing \citet{daugherty1991credibility} and naming the declining-nondeilning-yield phenomenon. The dynamic inconsistency reproduced here is the analytic root of that practical drift, developed further in the author's broader wood-supply-modelling research program \citep{paradis2017compiling,paradis2018retrofitting,paradis2018estimating}. That the underlying analysis has been effectively inaccessible for three decades---an archival thesis obtainable only as a roughly \$200 custom print---plausibly explains why the phenomenon keeps being rediscovered. A citable, reproducible reference is overdue.

\hypertarget{remedies}{%
\subsection{Remedies}\label{remedies}}

The thesis and the subsequent literature point to three families of remedies, which the reproduction makes concrete. First, \emph{consistent formulations}: rather than solving one open-loop problem, compute solutions that are optimal for the sequential game the planning process actually is---subgame-perfect or feedback solutions in which each period's decision is optimal given the state and the continuation policy. The thesis discusses the generation of consistent solutions; computing them exactly for the full case study is a natural follow-on enabled by this stack. Second, \emph{credible precommitment}: inconsistency arises because the future planner is free to re-solve; institutional commitments that genuinely bind future decisions---for example, regeneration mandates attached to harvest---change the future subproblem and can restore credibility to announced plans. (The reproduction's harvest actions regenerate to age zero by construction, which makes this channel visible in the model.) Third, \emph{receding-horizon practice with honest reporting}: if planning is in fact sequential replanning, then the simulator used here---not the single open-loop solve---is the honest model of what the planning process will deliver, and plan documents should report the realized replanned trajectory rather than the open-loop projection.

\hypertarget{limitations}{%
\subsection{Limitations}\label{limitations}}

Three limitations bound the claims. First, the data vintage: the case study is built from the 1990 FEIS (the FORPLAN model lineage with updated 1980s prices), cross-checked against the 1987 DEIS that Daugherty used directly; the two vintages' yield tables agree, and the thesis's exact base-case schedule is not matched period by period because the thesis's per-period realized schedule is not printed. Validation is against the thesis's printed anchors (Tables 5.3 and 5.4), which the assembled case study reproduces in sign, ordering, and magnitude range. Second, the formulation caveat: the thesis used Model II; the reproduction uses Model I (Section 2.2). Third, scope: this is one case study and one simulator design (rolling horizon, fixed seeds). The pervasiveness result (100\% of grid cells) is a property of this grid, and while it reproduces the thesis's qualitative conclusion over a wide range of conditions, extension to other forests, formulations, and consistent-solution constructs is future work---work that the open stack is intended to make easy.
